% The two studies every model has: fewer training games, and one group of
% hand features left out at a time. The same for every model, so a paper
% reads this after changes.tex and what it says of its own runs there, and
% adds what its own numbers mean. Each study waits for its runs.
\ifgames
\paragraph{Fewer training games.}
How much of a score is the data, and would more of it help? The games study
trains a run again on half and on a quarter of the training games, and
scores it on the same test rows (Table~\ref{tab:games},
Figure~\ref{fig:games}). The games are chosen by a hash, so the quarter lies
inside the half. A network is given twice the epochs on half the games and
four times on a quarter, so it takes as many steps each time and loses only
the games. The study is read forward: where halving the games costs little,
doubling them would buy little, so the curve says what a larger dataset is
worth before anyone builds one.

\begin{table}[htbp]\centering\footnotesize
\block{study-games}
\caption{The games study. Each run is trained on a share of the training
games and scored on all the test rows. Its gain is paired with the run that
was trained on all the games: below zero the run is worse, and bold when the
95\% interval leaves zero out.}
\label{tab:games}
\end{table}

\begin{figure}[htbp]\centering
\block{study-games-figure}
\caption{The games study as a learning curve: the score on the test rows by
the share of the training games a run was given. Each step to the right
doubles the games.}
\label{fig:games}
\end{figure}
\fi

\iffeatures
\paragraph{A group of hand features left out.}
The hand features come in groups: material, the moves on offer, pieces that
can be taken, passed pawns and attacks near the kings, and for a model that
reads a line of moves, its first move and the line as a whole. The features
study trains a run again without one group at a time
(Table~\ref{tab:features}, Figure~\ref{fig:features}). A run that is worse
without a group was using it, and nothing else the run reads holds the same.
A run that is no worse either made no use of the group or finds the same
elsewhere, in another group or on the board. So the study measures what a
group adds on top of everything else, which is less than what it knows: a
group can matter a great deal on a board and cost nothing here.

\begin{table}[htbp]\centering\footnotesize
\block{study-features}
\caption{The features study. Each run reads every hand feature but one
group. Its gain is paired with the run that reads them all: below zero the
run is worse without the group, and bold when the 95\% interval leaves zero
out. The group whose loss costs most comes first.}
\label{tab:features}
\end{table}

\begin{figure}[htbp]\centering
\block{study-features-figure}
\caption{What leaving a group of hand features out costs the reference: the
paired gain of the run without the group over the run with it, with its 95\%
interval. Left of the line the run is worse without the group.}
\label{fig:features}
\end{figure}
\fi
